% Generated draft: requires mathematical and source review.
\documentclass[11pt]{article}
\usepackage[margin=25mm]{geometry}
\usepackage{fontspec}
\setmainfont{Latin Modern Roman}
\usepackage{amsmath,amssymb,mathtools,unicode-math}
\setmathfont{Latin Modern Math}
\usepackage{enumitem,needspace,xurl,hyperref}
\setlength{\parindent}{0pt}
\setlength{\parskip}{5pt}
\setlength{\emergencystretch}{3em}
\newcommand{\N}{\mathbb{N}}
\newcommand{\Z}{\mathbb{Z}}
\newcommand{\Q}{\mathbb{Q}}
\newcommand{\R}{\mathbb{R}}
\newcommand{\C}{\mathbb{C}}
\newcommand{\F}{\mathbb{F}}
\newcommand{\A}{\mathbb{A}}
\newcommand{\PP}{\mathbb{P}}
\newcommand{\E}{\mathbb{E}}
\newcommand{\eps}{\varepsilon}
\newcommand{\dd}{\,\mathrm{d}}
\begin{document}
\section*{UnsolvedMath ID: 10400047}
\textbf{Source problem code:} catalogue entry
\subsection{Definitions and mathematical statement}
A virtual knot is represented by a planar knot diagram with ordinary over/under crossings and virtual crossings, modulo the generalized Reidemeister moves and detour equivalence. It can also be modeled by a knot in a thickened oriented surface, allowing the appropriate stabilization away from the knot. The diagram and surface conventions must be kept compatible.

The problem attributed to M. Polyak is to establish a Goussarov--Habiro theory for virtual knots. In the classical setting this terminology refers to clasper/surgery moves, the resulting finite-type filtration and equivalence relations, and their relationship with diagrammatic and finite-type invariant structures. The source does not state which of these conclusions or degree conventions are required in the virtual setting.

A precise extension must choose allowed virtual clasper or surgery operations, prove invariance under virtual equivalence (including surface stabilization if used), and identify the intended finite-type notion and associated graded objects. Virtual-knot finite-type theories based on different singularization/virtualization rules need not be the same, so the mere existence of some finite-type invariants does not determine or establish the requested theory. No classical degree-shift or classification theorem is transferred without specifying its virtual counterpart.

Thus the source is a theory-development programme, not a single asserted isomorphism, move theorem or universal classification. Its local moves, coefficient ring, knot category (for example long versus closed and framed versus unframed), filtration and success criteria are not supplied.
\subsection{Short English statement}
Develop a virtual-knot counterpart of Goussarov--Habiro clasper and finite-type theory, with well-defined moves, filtration and invariant comparison. The intended virtual finite-type and knot-category conventions are missing.
\subsection{Sources}
\begin{itemize}
\item Local UnsolvedMath registry, permanent ID 10400047, exact source code catalogue entry, Problem 2.25; source attribution M. Polyak, no proposer/year in metadata. Registry status partially solved, classification PARTIAL-PROGRESS; the registry title is truncated.
\item Ohtsuki, \emph{Problems on invariants of knots and 3-manifolds} (2002), Problem 2.25, PDF page 48, \url{https://msp.org/gtm/2002/04/gtm-2002-04-024s.pdf}. Supplied pdf-font-structure extraction accessed July 29, 2026; provenance distinguishes 231 bold numbered items. No PDF was accessed here.
\item The supplied summary reports virtual finite-type work but no established full theory, without exact references or a criterion for full. No specific later theorem is inferred.
\item Provenance flags current primary-literature status and redistribution terms NEEDS\_REVIEW.
\end{itemize}
\end{document}
